\documentclass[pdflatex,sn-mathphys-ay]{sn-jnl}


\usepackage{graphicx}%
\usepackage{multirow}%
\usepackage{amsmath,amssymb,amsfonts}%
\usepackage{amsthm}%
\usepackage{mathrsfs}%
\usepackage{xcolor}%
\usepackage{textcomp}%
\usepackage{manyfoot}%
\usepackage{booktabs}%
\usepackage{algorithm}%
\usepackage{algorithmicx}%
\usepackage{algpseudocode}%
\usepackage{listings}%
\usepackage{siunitx}
\sisetup{group-separator={,}, group-minimum-digits=4}
\usepackage{tabularx}
\usepackage{dcolumn}
\usepackage{xurl}

%%%%%=============================================================================%%%%
%%%%  Remarks: This template is provided to aid authors with the preparation
%%%%  of original research articles intended for submission to journals published 
%%%%  by Springer Nature. The guidance has been prepared in partnership with 
%%%%  production teams to conform to Springer Nature technical requirements. 
%%%%  Editorial and presentation requirements differ among journal portfolios and 
%%%%  research disciplines. You may find sections in this template are irrelevant 
%%%%  to your work and are empowered to omit any such section if allowed by the 
%%%%  journal you intend to submit to. The submission guidelines and policies 
%%%%  of the journal take precedence. A detailed User Manual is available in the 
%%%%  template package for technical guidance.
%%%%%=============================================================================%%%%

%% as per the requirement new theorem styles can be included as shown below
\theoremstyle{thmstyleone}%
\newtheorem{theorem}{Theorem}%  meant for continuous numbers
%%\newtheorem{theorem}{Theorem}[section]% meant for sectionwise numbers
%% optional argument [theorem] produces theorem numbering sequence instead of independent numbers for Proposition
\newtheorem{proposition}[theorem]{Proposition}% 
%%\newtheorem{proposition}{Proposition}% to get separate numbers for theorem and proposition etc.

\theoremstyle{thmstyletwo}%
\newtheorem{example}{Example}%
\newtheorem{remark}{Remark}%

\theoremstyle{thmstylethree}%
\newtheorem{definition}{Definition}%

\raggedbottom
%%\unnumbered% uncomment this for unnumbered level heads

\begin{document}

\title[d]{d}



\author[1]{\fnm{Seungmin} \sur{Yoon}}\email{smyoon@ks.ac.kr}

\author[1]{\fnm{Sanghun} \sur{Kim}}\email{fatherofdamin@ks.ac.kr}

\author*[1]{\fnm{Hyunpyo} \sur{Jeon}}\email{hpjeon@ks.ac.kr}
%\equalcont{These authors contributed equally to this work.}

\affil[1]{Department of Biosafety, Kyungsung University, Busan, Republic of Korea}

\abstract{Write later}
%%================================%%
%% Sample for structured abstract %%
%%================================%%

% \abstract{\textbf{Purpose:} The abstract serves both as a general introduction to the topic and as a brief, non-technical summary of the main results and their implications. The abstract must not include subheadings (unless expressly permitted in the journal's Instructions to Authors), equations or citations. As a guide the abstract should not exceed 200 words. Most journals do not set a hard limit however authors are advised to check the author instructions for the journal they are submitting to.
% 
% \textbf{Methods:} The abstract serves both as a general introduction to the topic and as a brief, non-technical summary of the main results and their implications. The abstract must not include subheadings (unless expressly permitted in the journal's Instructions to Authors), equations or citations. As a guide the abstract should not exceed 200 words. Most journals do not set a hard limit however authors are advised to check the author instructions for the journal they are submitting to.
% 
% \textbf{Results:} The abstract serves both as a general introduction to the topic and as a brief, non-technical summary of the main results and their implications. The abstract must not include subheadings (unless expressly permitted in the journal's Instructions to Authors), equations or citations. As a guide the abstract should not exceed 200 words. Most journals do not set a hard limit however authors are advised to check the author instructions for the journal they are submitting to.
% 
% \textbf{Conclusion:} The abstract serves both as a general introduction to the topic and as a brief, non-technical summary of the main results and their implications. The abstract must not include subheadings (unless expressly permitted in the journal's Instructions to Authors), equations or citations. As a guide the abstract should not exceed 200 words. Most journals do not set a hard limit however authors are advised to check the author instructions for the journal they are submitting to.}

\keywords{K-REACH; genotoxicity; in silico; QSAR; structural alerts; machine learning}



\maketitle

\clearpage
\section{Introduction}\label{sec:intro}

Genotoxicity is a critical toxicological endpoint that assesses the potential of chemicals to induce DNA damage, 
mutations, or carcinogenesis and plays a central role in regulatory chemical evaluation \citep{Karamertzanis2025, Hayashi2022, EFSA2018}.



%규제 체계에서 시험을 의무화함에 따라 유전독성 데이터의 활용 가능성이 많아졌고, 이는 In silico model 개발을 촉진했다\citep{HASSELGREN2019104403}.

%현재 범용적인 QSAR 모델인 OECD QSAR Toolbox와 VEGA-QSAR는 산업용 화학물질에 적용 도메인이 알맞지 않음\citep{YOON2026100405}
\section{Materials and Methods}\label{sec:methods}

%\subsection{Data source}\label{subsec:data}

%\subsection{Data processing}\label{subsec:processing}

%\subsection{Modeling}\label{subsec:modeling}

\section{Results}\label{sec:results}


\section{Discussion}\label{sec:discussion}


\section{Limitations and conclusions}\label{sec:limitations}
 



\clearpage

\bibliography{sn-bibliography}
\clearpage
%\input{7. Supplementery 1}

\end{document}
